\documentclass[aspectratio=169, table, xcdraw]{beamer}

\usetheme{Madrid}
\usecolortheme{default}
\definecolor{ucbblue}{RGB}{0, 51, 102}

\setbeamercolor{structure}{fg=ucbblue}
\setbeamercolor*{palette primary}{fg=white,bg=ucbblue}
\setbeamercolor*{palette secondary}{fg=white,bg=ucbblue!80}
\setbeamercolor*{palette tertiary}{fg=white,bg=ucbblue!90}
\setbeamercolor{title}{fg=white, bg=ucbblue}
\setbeamercolor{frametitle}{fg=white, bg=ucbblue}
\setbeamercolor{item}{fg=ucbblue}

\usepackage[T1]{fontenc}
\usepackage[utf8]{inputenc}
\usepackage{tgpagella}
\usefonttheme{serif}
\usepackage[spanish, es-noshorthands]{babel}
\usepackage{amsmath, amssymb, bm}
\usepackage{graphicx}
\usepackage[most]{tcolorbox}
\usepackage{booktabs} % Soluciona los errores \toprule, \midrule, \bottomrule
\usepackage{tabularx} % Soluciona el Overfull \hbox permitiendo saltos de línea
\usepackage{tikz}
\usepackage{pgfplots}
\pgfplotsset{compat=1.18}
\usepackage[american]{circuitikz}
\usetikzlibrary{shapes.geometric, arrows.meta, positioning, calc}

\title[BJT y Polarización]{El Transistor BJT: Polarización DC y Recta de Carga}
\subtitle{IMT-121 Circuitos Electrónicos I — Unidad 2}
\author[B. Quiroga Turdera]{M.Sc. Ing. Bernardo Quiroga Turdera}
\institute[UCB Tarija]{Universidad Católica Boliviana ``San Pablo'' — Sede Tarija}
\date{Semana 6}

\begin{document}

\begin{frame}
    \titlepage
\end{frame}

\begin{frame}{Hilos Conductores: De la Unidad 1 a la Unidad 2}
    \begin{block}{Evolución del Análisis}
        \centering
        Protección $\rightarrow$ Señal pequeña $\rightarrow$ \textbf{Transistor Activo} $\rightarrow$ \textbf{Punto Q estable} $\rightarrow$ Par Diferencial
    \end{block}
    
    \vspace{0.3cm}
    \textbf{Pregunta Guía de Ingeniería:}
    \begin{center}
        \textit{¿Cómo situamos un transistor en un punto de operación predecible (DC) para que una pequeña señal de un sensor pueda ser amplificada posteriormente sin recortes ni pérdidas?}
    \end{center}
\end{frame}

\begin{frame}{Estructura Física del BJT (NPN)}
    \begin{columns}
        \begin{column}{0.48\textwidth}
            \textbf{Regiones Semiconductoras:}
            \begin{itemize}
                \item \textbf{Emisor (E):} inyecta portadores.
                \item \textbf{Base (B):} región delgada de control.
                \item \textbf{Colector (C):} recolecta portadores.
            \end{itemize}
            
            \vspace{0.2cm}
            \textbf{Ley de Corrientes de Kirchhoff (KCL):}
            \begin{equation*}
                I_E = I_C + I_B
            \end{equation*}
        \end{column}
        \begin{column}{0.48\textwidth}
            \centering
            \begin{circuitikz}[scale=1.2, transform shape, thick]
                \draw (0,0) node[npn, tr circle](Q){};
                \node [right] at (Q.C) {C ($I_C$)};
                \node [right] at (Q.E) {E ($I_E$)};
                \node [left] at (Q.B) {B ($I_B$)};
                \draw[->, blue] (Q.C) ++(0,0.5) -- ++(0,-0.4);
                \draw[->, blue] (Q.B) ++(-0.5,0) -- ++(0.4,0);
                \draw[->, blue] (Q.E) ++(0,-0.2) -- ++(0,-0.4);
            \end{circuitikz}
        \end{column}
    \end{columns}
\end{frame}

\begin{frame}{Regiones de Operación}
    El modelo del circuito depende de la región de operación. 
    \textbf{Aproximación para análisis manual:} $V_{BE} \approx 0.7\,\text{V}$ (solo si está encendido).
    
    \vspace{0.2cm}
    \begin{table}[]
        \centering
        \small
        \renewcommand{\arraystretch}{1.3}
        \begin{tabularx}{\textwidth}{|l|X|X|X|}
        \hline
        \rowcolor{ucbblue!20} 
        \textbf{Región} & \textbf{Unión BE} & \textbf{Unión BC} & \textbf{Interpretación en el Circuito} \\ \hline
        Corte & No polarizada en directa & No polarizada en directa & Transistor aprox. apagado \\ \hline
        Activa Directa & Polarizada en directa & Polarizada en inversa & Región de \textbf{amplificación} \\ \hline
        Saturación & Polarizada en directa & Polarizada en directa & Comportamiento de interruptor \\ \hline
        \end{tabularx}
    \end{table}

    \begin{alertblock}{¡Cuidado!}
        La relación $I_C \approx \beta I_B$ \textbf{no} es una ley universal. Solo es válida en la región Activa Directa, que requiere: $V_{BC} < 0$ (o $V_{CE} > V_{BE}$).
    \end{alertblock}
\end{frame}

\begin{frame}{Circuito de Polarización DC y Método de Verificación}
    \begin{columns}
        \begin{column}{0.4\textwidth}
            \centering
            \begin{circuitikz}[scale=0.9, transform shape, thick]
                \draw (0,0) node[npn](Q){};
                \draw (Q.E) node[ground]{};
                \draw (Q.C) to[R, l_=$R_C$, i_<=$I_C$] (0,3);
                \draw (Q.B) -- (-2.5, 0 |- Q.B) to[R, l=$R_B$, i_<=$I_B$] (-2.5, 3);
                \draw (-2.5,3) -- (0,3) -- (0,3.5) node[above]{$+V_{CC}$};
                
                \draw[->, red] (-0.8, -0.3) arc (-180:-90:0.5) node[midway, left]{$V_{BE}$};
                \draw[->, red] (1, 2) arc (90:-90:1) node[midway, right]{$V_{CE}$};
            \end{circuitikz}
        \end{column}
        \begin{column}{0.55\textwidth}
            \textbf{Método Lógico:}
            \begin{enumerate}
                \item \textbf{Asumir} modelo (Activa: $V_{BE}=0.7\,\text{V}$, $I_C = \beta I_B$).
                \item \textbf{Resolver} la red DC.
                \item \textbf{Verificar} la suposición ($V_{CE} > 0.7\,\text{V}$).
                \item \textbf{Aceptar o Rechazar} el modelo.
            \end{enumerate}
            
            \vspace{0.2cm}
            \textbf{Malla Base:} \quad $I_B = \frac{V_{CC} - V_{BE}}{R_B}$ \\
            \textbf{Malla Colector:} \quad $V_{CE} = V_{CC} - I_C R_C$
        \end{column}
    \end{columns}
\end{frame}

\begin{frame}{Ejemplo Guiado: Análisis DC Completo}
    \textbf{Datos:} $V_{CC}=10\,\text{V}$, $R_B=470\,\text{k}\Omega$, $R_C=2.0\,\text{k}\Omega$, $\beta=100$, $V_{BE}=0.7\,\text{V}$.
    
    \vspace{0.2cm}
    1. Corriente de Base ($I_B$):
    \begin{equation*}
        I_B = \frac{10\,\text{V} - 0.7\,\text{V}}{470\,\text{k}\Omega} = 19.79\,\mu\text{A}
    \end{equation*}
    
    2. Candidato a Corriente de Colector ($I_C$):
    \begin{equation*}
        I_C = \beta I_B = 100(19.79\,\mu\text{A}) = 1.979\,\text{mA}
    \end{equation*}
    
    3. Verificación de Tensión ($V_{CE}$ y $V_{BC}$):
    \begin{equation*}
        V_{CE} = 10\,\text{V} - (1.979\,\text{mA})(2\,\text{k}\Omega) = 6.04\,\text{V}
    \end{equation*}
    \begin{equation*}
        V_{BC} = V_{BE} - V_{CE} = 0.7\,\text{V} - 6.04\,\text{V} = -5.34\,\text{V} < 0 \quad \mathbf{(\text{¡Consistente!})}
    \end{equation*}
    El Punto Q es: $\mathbf{Q \approx (1.98\,\text{mA},\, 6.04\,\text{V})}$.
\end{frame}

\begin{frame}{Ejemplo: Cuando la Suposición Falla}
    Mismo circuito, pero reducimos fuertemente $R_B$: $\mathbf{R_B = 100\,\text{k}\Omega}$.
    
    \vspace{0.2cm}
    1. Corriente de Base ($I_B$):
    \begin{equation*}
        I_B = \frac{10\,\text{V} - 0.7\,\text{V}}{100\,\text{k}\Omega} = 93\,\mu\text{A}
    \end{equation*}
    
    2. Candidato a Corriente de Colector ($I_C$):
    \begin{equation*}
        I_C = \beta I_B = 100(93\,\mu\text{A}) = 9.3\,\text{mA}
    \end{equation*}
    
    3. Verificación de Tensión ($V_{CE}$):
    \begin{equation*}
        V_{CE} = 10\,\text{V} - (9.3\,\text{mA})(2\,\text{k}\Omega) = \mathbf{-8.6\,\text{V}}
    \end{equation*}
    
    \begin{alertblock}{Conclusión Correcta}
        Un $V_{CE}$ negativo es \textbf{físicamente inconsistente} para este circuito. \textbf{Rechazamos} el modelo activo; el transistor está en (o fuertemente impulsado hacia) \textbf{saturación}.
    \end{alertblock}
\end{frame}

\begin{frame}{La Recta de Carga DC y el Punto Q}
    \begin{columns}
        \begin{column}{0.48\textwidth}
            La recta de carga es impuesta por el \textbf{circuito externo} al colector.
            \begin{equation*}
                V_{CE} = V_{CC} - I_C R_C \implies I_C = \frac{V_{CC} - V_{CE}}{R_C}
            \end{equation*}
            
            \textbf{Intersecciones:}
            \begin{itemize}
                \item Eje $V_{CE}$ ($I_C=0$): $V_{CE} = V_{CC}$
                \item Eje $I_C$ ($V_{CE}=0$): $I_C = \frac{V_{CC}}{R_C}$
            \end{itemize}
            El Punto Q reside sobre esta recta de carga.
        \end{column}
        \begin{column}{0.48\textwidth}
            \centering
            \begin{tikzpicture}
                \begin{axis}[
                    axis lines=left,
                    xlabel={$V_{CE}$ (V)},
                    ylabel={$I_C$ (mA)},
                    xmin=0, xmax=11,
                    ymin=0, ymax=6,
                    grid=both,
                    width=7cm, height=5.5cm,
                    xtick={0,2,4,6,8,10},
                    ytick={0,1,2,3,4,5}
                ]
                % Load line for Vcc=10, Rc=2k -> max Ic=5mA
                \addplot[thick, ucbblue, domain=0:10] {(10-x)/2};
                \addplot[mark=*, red, mark size=2.5pt] coordinates {(6.04, 1.979)} node[above right, font=\bfseries]{Q};
                \node at (axis cs: 10, 0) [above left] {Corte};
                \node at (axis cs: 0, 5) [right] {Saturación};
                \end{axis}
            \end{tikzpicture}
        \end{column}
    \end{columns}
\end{frame}

\begin{frame}{Sensibilidad a $\beta$ y Estabilidad de Polarización}
    En polarización fija, $I_C \propto \beta$. Dado que $\beta$ varía drásticamente (por fabricación o temperatura), este circuito es frágil. \\
    Ejemplo con $R_B=470\,\text{k}\Omega$, $R_C=2\,\text{k}\Omega$:
    
    \vspace{0.3cm}
    \centering
    \renewcommand{\arraystretch}{1.3}
    \begin{tabular}{ccc}
    \toprule
    \textbf{Parámetro $\beta$} & \textbf{$I_C$ Estimada} & \textbf{$V_{CE}$ Resultante} \\ \midrule
    50 & 0.989\,\text{mA} & 8.02\,\text{V} \\
    100 & 1.979\,\text{mA} & 6.04\,\text{V} \\
    200 & 3.957\,\text{mA} & 2.09\,\text{V} \\ \bottomrule
    \end{tabular}
    
    \vspace{0.3cm}
    \begin{block}{Motivación para Diseños Futuros}
        Una polarización robusta buscará independizar $I_C$ de variaciones en $\beta$.
    \end{block}
\end{frame}

\begin{frame}{Puente hacia el Análisis de Pequeña Señal}
    El modelo de pequeña señal no reemplaza al modelo DC, \textbf{depende de él}.
    
    Las señales en el circuito son perturbaciones locales alrededor del Punto Q:
    \begin{align*}
        i_C(t) &= I_{CQ} + i_c(t) \\
        v_{BE}(t) &= V_{BEQ} + v_{be}(t) \\
        v_{CE}(t) &= V_{CEQ} + v_{ce}(t)
    \end{align*}
    
    \vspace{0.2cm}
    \textbf{Parámetros dinámicos:} La transconductancia ($g_m$) o resistencia ($r_\pi$) no pueden ser elegidos sin primero establecer $Q$:
    \begin{equation*}
        g_m = \frac{I_{CQ}}{V_T}
    \end{equation*}
\end{frame}

\end{document}
